\sectionguard
\section*{The Propers: Interpretive Possibilities}

\begin{studybox}{What this section is}
Everything below is exploratory editorial and AI proposal. None of it is
sourced historical intent, attributed teaching, or documented reception, and
none of it should be quoted as any of those. Each proposal was tested against
the corpus checked for this guide and against a targeted search for its own
distinctive conjunction; the results of those searches, and each proposal's
controlling limit, are recorded in \texttt{research/scope.md}. Evidence,
doctrine, and the literal sense of the appointed texts govern; where a
proposal and a checked witness conflict, the witness wins.
\end{studybox}

\begin{proposal}{The Gradual corrects the Introit's verb, and the Epistle explains why it can}
\pfield{Anchors}{Introit \cue{Int.}, \latin{ne derelínquas in finem};
Gradual \cue{Grad.}, \latin{ne obliviscáris in finem}; Epistle \cue{Ep.},
\latin{testaméntum confirmátum a Deo \dots\ non írritum facit}.}
\pfield{Mechanism}{The two chants sing the same verse with different verbs:
the Introit pleads against being \emph{forsaken}, the Gradual against being
\emph{forgotten} --- and the Gradual's verb is the psalter's own. Between
them stands a lesson about a \latin{testaméntum}: a covenant reduced to an
instrument, confirmed, and therefore --- as Ambrosiaster and Aquinas develop
the image --- as unforgettable as a will in probate. The day's movement runs
from the fear of abandonment to the fear of oblivion to the answer that
covers both: what is confirmed in writing can be neither dropped nor
forgotten.}
\pfield{Fruit}{The two petitions stop being synonyms. Abandonment is answered
by presence, oblivion by the covenant's permanence; the formulary supplies a
different consolation for each fear.}
\pfield{What the ordinary reading misses}{Element by element, the verbal
variation between Introit and Gradual reads as chant-transmission noise, and
most hand missals translate the two verbs identically.}
\pfield{Strongest limit}{The variation almost certainly \emph{is} Old-Latin
psalter transmission --- no checked witness comments on it, and nothing
documents a compiler choosing one verb against the other. The proposal reads
the printed page, not any intention behind it.}
\end{proposal}

\begin{proposal}{The voices God is asked not to forget are the ten that cried and the one that came back}
\pfield{Anchors}{Introit \cue{Int.}, \latin{ne obliviscáris voces
quæréntium te}; Gospel \cue{Gosp.}, \latin{levavérunt vocem \dots\ cum magna
voce magníficans Deum}.}
\pfield{Mechanism}{Where the Clementine asks God to remember the voices of his
\emph{enemies}, the chant sings the voices of his \emph{seekers}. The Gospel
then stages exactly two voices of seekers: ten lifted up together for mercy,
one returned loudly in thanks --- and the nine's cleansing stands while their
voice never comes back. Read together, the chant's plea and the narrative
divide one prayer into its two halves, petition and return, and show which half
goes missing.}
\pfield{Fruit}{Thanksgiving stops being politeness after the fact and becomes
the second half of the voice that asked --- the half by which a seeker stays
a seeker after being answered.}
\pfield{What the ordinary reading misses}{The chant's ``seekers'' is invisible
in every English a reader will hold, since the Douay renders the Clementine's
``enemies''; without it the chant has no voice-vocabulary to share with the
Gospel.}
\pfield{Strongest limit}{The proposal must not be built on a supposed chant
liberty, because there is none: the commentary shows that the seekers are the
Septuagint's reading and the enemies the Hebrew's, that Augustine, Cassiodorus
and Theodoret all had the seekers in front of them, and that Migne's apparatus
sets the whole textual situation out. So the wording carries no compiler's
intention at all --- it is simply the psalter the Roman chant inherited. Beyond
that, \latin{vox} is among the commonest words in the psalter and the Gospels,
and no checked witness joins these two texts.}
\end{proposal}

\begin{proposal}{Three clocks: the generations, the four hundred thirty years, and my times}
\pfield{Anchors}{Alleluia \cue{All.}, \latin{refúgium \dots\ a generatióne et
progénie}; Epistle \cue{Ep.}, \latin{post quadringéntos et trigínta annos};
Offertory \cue{Off.}, \latin{in mánibus tuis témpora mea}.}
\pfield{Mechanism}{The formulary counts time three ways. The Alleluia counts
in generations, and finds the refuge constant across all of them --- in the
one psalm bearing Moses' name, whose burden is the contrast between God's
eternity and our brevity. The Epistle counts in years, and finds that four
hundred and thirty of them change nothing that God has confirmed. The
Offertory stops counting and hands the clock over: \latin{témpora mea} ---
the old psalter's ``times'' where the Clementine has ``lots,'' time standing
where chance stood. Three scales, one claim: duration cannot erode what God
holds.}
\pfield{Fruit}{The Introit's anguished \latin{Ut quid \dots\ in finem?} ---
how long, why to the end --- receives an answer about time itself rather
than about the crisis of the moment: delay is not abandonment, because the
interval belongs to the same hands as the promise.}
\pfield{What the ordinary reading misses}{Taken singly, the Alleluia is a
generic comfort, the 430 years a Pauline technicality, and the Offertory an
ejaculation of trust; the day's preoccupation with measured time only appears
when the three are read as one series.}
\pfield{Strongest limit}{No checked witness connects any two of these three
texts through time-vocabulary, and the nearest thing to an analogue cuts only
one way: Augustine, expounding the Alleluia's verse, makes the refuge God's
\latin{aeternitas} and what one flees to it from \latin{temporis mutabilitas},
and Theodoret glosses the Offertory's variant as the changes of circumstance
--- but neither joins the two chants, and neither mentions the four hundred and
thirty years. \latin{Témpora}, \latin{anni} and \latin{generatióne} share no
root; the Offertory's reading predates any pairing with this Epistle; and chant
assignments are commonly musical and cyclical, not thematic.}
\end{proposal}

\begin{proposal}{The keeper the chants pray for appears in the Gospel as the man who came back}
\pfield{Anchors}{Introit and Gradual \cue{Int., Grad.}, the pleas
\latin{réspice \dots\ ne derelínquas \dots\ ne obliviscáris}; Gospel
\cue{Gosp.}, \latin{et hic erat Samaritánus}, with Augustine's documented
etymology \latin{Samaritanus = custos}, ``keeper.''}
\pfield{Mechanism}{The chants beg God to keep --- to look upon, not forsake,
not forget. The Gospel's one returning figure is named by a word the
tradition heard as ``keeper'': Augustine derives it so and lets the leper keep
the unity of the kingdom by giving back what he received; Rupert of Deutz, as
Guéranger reports him, read this whole Sunday's office through
\latin{Samaritanus} as \latin{custos}. One further fact tightens the join and
is documented in the commentary: Augustine makes the same etymology a second
time, in his exposition of the \emph{Offertory's} psalm, at the verse that asks
God to be a house of refuge --- ``Samarites enim custos interpretatur \dots\
sic intellegi volens nostrum se esse custodem.'' The proposal is the join: the
day that pleads for keeping is given, as its narrative center, a keeper ---
covenant faithfulness answered not only from God's side but modeled from the
creature's, in the single act of returning thanks.}
\pfield{Fruit}{Gratitude is reclassified: not an ornament after grace but the
creature's half of covenant-keeping, the act by which what was given is held
rather than merely received.}
\pfield{What the ordinary reading misses}{The etymology is invisible in
translation, and the office-wide \latin{custos} reading rode on the Good
Samaritan Gospel that this Sunday carried in an older arrangement --- so the
modern page shows no trace of either.}
\pfield{Strongest limit}{Augustine applies \latin{custos} to this leper, but
Rupert applied it to the \emph{Good} Samaritan as Christ, a different
pericope assigned to this Sunday only in a pre-Tridentine ordering; the
transfer of that office-wide reading onto the 1962 chants is this proposal's
own act. Hebrew etymology of \latin{Samaritanus} as ``keeper'' is itself a
patristic reception, not a modern lexical judgment.}
\end{proposal}

\begin{proposal}{The two askings of \latin{augméntum} bracket the nine who stopped}
\pfield{Anchors}{Collect \cue{Coll.}, \latin{fidei, spei et caritátis
augméntum}; Postcommunion \cue{Postcomm.}, \latin{redemptiónis ætérnæ \dots\
augméntum}; Gospel \cue{Gosp.}, \latin{nonne decem mundáti sunt? et novem ubi
sunt?}}
\pfield{Mechanism}{The day's first prayer and its last share one noun:
increase. Between them stands a Gospel in which ten receive a grace entire
--- they are cleansed, nothing lacking --- and nine go no further, while the
one who returns is moved from cleansing to salvation. Bede states the
distinction: faith's working, not its bare knowledge, saves; Trent quotes
this very Collect for the increase of justification. The proposal reads the
bracket as the formulary's own diagnosis of the nine: what they lacked was
not the gift but the increase --- the movement that turns a grace received
into a grace grown.}
\pfield{Fruit}{The frightening arithmetic of the Gospel (one in ten) becomes
usable: the difference between the nine and the one is presented as
something the day teaches you to pray for, twice.}
\pfield{What the ordinary reading misses}{The bracket is invisible in the
registered English, which renders the first \latin{augméntum} ``increase''
but drops the second entirely --- only the Latin shows the day ending with
the word it began with.}
\pfield{Strongest limit}{The Collect and Postcommunion were not composed for
this Gospel --- the pair already travels together in the Gelasian
sacramentary's unnumbered Sunday masses, centuries before any fixed pericope
--- so the bracket is a fact of the printed formulary, not of anyone's
design; and Bede's distinction concerns the Gospel alone, not the orations.}
\end{proposal}
